\documentclass[12pt]{article}
\usepackage{thianswersheetBB}

\usepackage{ifxetex}
\ifxetex
\usepackage{xltxtra} 
\setmainfont[
Path=Charis Font/,
BoldFont       = CharisSILB.ttf,
ItalicFont     = CharisSILI.ttf,
BoldItalicFont = CharisSILBI.ttf
]{CharisSILR.ttf}
\else
\usepackage[utf8]{vntex}
\usepackage{lmodern}
\usepackage{charter}
\fi
\usepackage[
marginparsep=6pt,
top=1.5cm,
bottom=2cm,
left=1cm,
right=1cm,
bindingoffset=0.5cm,
footskip=0.75cm,
twoside,
a4paper]{geometry}
\usepackage{fancyhdr}
\pagestyle{fancy}
\fancyhf{}
\renewcommand{\headrulewidth}{0pt}
\tentruonghoc{ĐH DUY TÂN}
\tendonvi{Khoa: MT \& KHTN}
\tenkithi{KÌ THI KẾT THÚC HỌC PHẦN}
\tendethiso{1}
\dauphantachgiuaSodevaMade{.}
\soluongphienbanmadethi{4} %4 là số lượng mã đề thi khác nhau được in ra
\chuphienbanmadethi{\Alph}%{\alph}%{\Roman}

% THIẾT LẬP TỔNG ĐIỂM CỦA CÁC PHẦN
\tongdiemcauhoiTracNghiem{4.0}
\tongdiemcauhoiNgan{3.0}
\tongdiemcauhoiTuLuan{3.0}

% THIẾT LẬP CÁC CÂU HỎI TNKQ
\socauhoiTracNghiem{20}
\sophuongantraloiTracNghiem{4}
\socauhoiTracNghiemtrenDong{20}% chia số lượng câu hỏi Trắc Nghiệm trong 1 dòng

% THIẾT LẬP CÁC CÂU HỎI NGẮN
\socauhoiNgan{6}  % số lượng câu hỏi ngắn
\sodongmoicauhoiNgan{3}


\mauchucauhoi{red} % chọn màu cho chữ câu hỏi
\tenhocphan{Lí thuyết xác suất \& thống kê toán}
\tenmamon{STA151}
\thoigianlambai{75}

\begin{document}

\ingiaylambaiCauHoiDongBB{true}{blue}% số câu hỏi được xếp theo Cột, bạn có thể thay true bởi false để xem hiển thị
%true: hiển thị mã đề trên mỗi giấy làm bài!
%false: KHÔNG thiển thị mã đề trên giấy làm bài

\phantieudeTrangAnswerBB

\answerCauHoiDongBB
{4;3;2;3;1;3;2;3;1;2;0;2;2;1;3;2;2;1;2;4}{blue}{1}%
%0: ko chọn đáp án ở câu này; 1,2,3,... đáp án tương ướng là A, B, C,...
%{1}: mã đề A

\answerCauHoiDongBB {4;3;2;3;2;4;3;2;3;1;2;0;1;2;2;1;3;2;2;1}{blue}{2}%

\answerCauHoiDongBB {4;3;2;3;2;4;3;2;3;1;2;0;1;2;2;1;3;2;2;1}{blue}{3}%

\answerCauHoiDongBB {4;3;2;3;1;2;2;1;3;2;2;1;2;4;3;2;3;1;2;0}{blue}{4}%
\end{document}



